\section{Related Works}
This section reviews and describes related works. The increasing complexity and volume of scientific literature make it challenging for researchers to conduct comprehensive, coherent, and well-structured reviews of the state of the art. In this context, 'Design, development, and implementation of an intelligent system based on large language model (LLM)-driven agents for the automation of academic section writing in scientific articles' aims to contribute by incorporating an agent-guided approach that makes autonomous intermediate decisions and provides comprehensive automation of the critical synthesis process, thereby reducing researchers' manual effort while enhancing the consistency and depth of analysis.

\par\par
Previous literature has been organized into three main categories: LLM Intelligent Agents, LLM Multi-Agent Systems, and LLM Human-Agent Collaboration. The 'LLM Intelligent Agents' category groups studies on the definitions, methodologies, applications, and challenges of intelligent agents based on large language models. The 'LLM Multi-Agent Systems' category focuses on research into architectures, progress, and challenges associated with systems composed of multiple LLM-driven agents. Finally, the 'LLM Human-Agent Collaboration' category examines the interaction and collaboration between humans and intelligent agents based on large language models, including their systems and applications.

\subsection{LLM Intelligent Agents}
The work 'EXPLORING LARGE LANGUAGE MODEL BASED INTELLIGENT AGENTS: D EFINITIONS , METHODS , AND PROSPECTS' \cite{ref-1} addresses research on intelligent agents based on large language models, focusing on their definitions, methods, and prospects. The methodology employed consists of a comprehensive review of the scientific literature on the subject. Among its key contributions, the study highlights the relevance of these agents for applications in artificial intelligence, robotics, and decision-making, while analyzing their current limitations and proposing future perspectives to improve their performance and generalizability. However, critical limitations identified include dependence on large volumes of data, the inherent complexity of machine learning, and deficiency in knowledge transfer between different tasks, which underscores the need for additional research to overcome these obstacles.

\par\par
The work 'Optimización de Agentes basados en Lenguaje Grande' \cite{ref-3} (Optimization of Large Language Model-based Agents) addresses the need to investigate the optimization of Large Language Model (LLM)-based agents to improve their performance in complex tasks, such as long-term planning and dynamic interaction with environments. The methodology focuses on optimizing LLM agents through fine-tuning techniques, reinforcement learning (RL), and other strategies, classifying methods into parameter-based fine-tuning, parameter-based RL, and hybrid optimization. As a key contribution, the study presents a systematic review of LLM agent optimization, identifying current limitations and opportunities for future development, in addition to providing a detailed analysis of the mentioned techniques. Nevertheless, the critical limitations of LLM agents reside in their long-term planning, dynamic interaction with environments, and adaptability to new scenarios, making the optimization of these agents essential to overcome these barriers and improve their performance in complex tasks.

\par\par
The work 'Large Language Model Agent: A Survey on Methodology, Applications and Challenges' \cite{ref-4} addresses the problem of the constant evolution of research on large language model (LLM) agents, highlighting the need for a systematic analysis to understand their architectures, collaboration, and evolution. The methodology focuses on creating a methodological taxonomy that deconstructs LLM agent systems into their fundamental components, including profile definition, memory mechanics, planning capabilities, and action execution. As a key contribution, the research presents a methodological taxonomy that connects the construction, collaboration, and evolution of LLM agent systems, also examining tools, communication protocols, and applications in this field. However, the critical limitation of the study is that, by focusing on the creation of a methodological taxonomy, its scope might be restricted to other aspects of LLM agent research.

\subsection{LLM Multi-Agent Systems}
The work 'Large Language Model based Multi-Agents: A Survey of Progress and Challenges' \cite{ref-2} (2023) addresses the problem that, while research on multi-agent systems based on large language models has shown significant progress, it also faces important challenges for its effective development and application. The methodology consists of a systematic review of research on these systems, highlighting their progress and challenges, and focusing on how they work, how agents are profiled, how they communicate, and how they acquire capabilities. As a key contribution, the study presents a general architecture for multi-agent systems based on large language models, emphasizing the importance of the environment-entity interface, agent profiling, inter-agent communication, and capability acquisition, in addition to a detailed analysis of their current applications. Nevertheless, critical limitations include the complexity of environments and the need to improve the efficiency and scalability of these systems.

\subsection{LLM Human-Agent Collaboration}
The work 'LLM-Based Human-Agent Collaboration and Interaction Systems: A Survey' \cite{ref-5} (2025) addresses the significant lack of research on collaboration systems between humans and LLM-based agents (LLM-Based Human-Agent Systems) in the field of artificial intelligence. The methodology employed consists of a comprehensive review of these systems, covering their key components, applications, challenges, and opportunities. As a key contribution, the study provides a complete and structured overview of human-agent collaboration systems based on LLMs, which is expected to foster research and development in this field. However, the main limitation of the paper is the lack of depth in some specific aspects of human-agent collaboration systems based on LLMs.

\par\par
Table \ref{tab:related_works_comparison} presents a comparison between the analyzed works and the approach proposed in this article. The table is based on specific criteria, namely:
\begin{itemize}
    \item \textbf{Title}: represents the main problem of the work or the proposal.
    \item \textbf{Main Problem Addressed}: describes the challenge or problem that each work attempts to address.
    \item \textbf{Key Limitation Reported}: identifies the key limitations or restrictions reported by the authors of each work.
    \item \textbf{Agent-LLM Integration Challenge}: analyzes specific challenges related to the integration of agents and large language models (LLMs).
    \item \textbf{Risk of Academic Rigor/Coherence}: evaluates the potential risk of losing academic rigor or coherence in the generated content.
    \item \textbf{Data/Base Model Dependency}: highlights the dependency on the base model and the data required for each work.
    \item \textbf{Implication for Academic Automation}: analyzes the implications of these critical points for the objective of automating academic tasks.
\end{itemize}

\begin{table}[h]
    \centering
    \caption{Comparative Analysis of Related Works and the Proposed Approach}
    \label{tab:related_works_comparison}
    \footnotesize % Adjust font size to fit table
    \begin{tabular}{|p{3.5cm}|p{3.5cm}|p{3.5cm}|p{3.5cm}|p{3.5cm}|p{3.5cm}|p{3.5cm}|}
        \hline
        \textbf{Title} & \textbf{Main Problem Addressed} & \textbf{Key Limitation Reported} & \textbf{Agent-LLM Integration Challenge} & \textbf{Risk of Academic Rigor/Coherence} & \textbf{Data/Base Model Dependency} & \textbf{Implication for Academic Automation} \\
        \hline
        Design, development, and implementation of an intelligent system based on large language model (LLMs)-driven agents for the automation of academic section writing in scientific articles & Increasing complexity, volume of scientific literature, comprehensive, coherent, and structured reviews. & Structural heterogeneity of PDFs, dependence on base model capabilities, information disambiguation. & Agents making autonomous decisions, critical synthesis, academic rigor. & Loss of academic rigor, disambiguating incomplete, implicit information. & Dependence on base model capabilities, information disambiguation. & Comprehensive automation, reduces manual effort, improves consistency, depth. \\
        \hline
        EXPLORING LARGE LANGUAGE MODEL BASED INTELLIGENT AGENTS: D EFINITIONS , METHODS , AND PROSPECTS \cite{ref-1} & Research on intelligent LLM agents, definitions, methods, prospects. & Need for large data, complexity of machine learning, lack of knowledge transfer. & Improve performance, generalizability, overcome current limitations. & Not specified & Need for large amounts of data. & Not specified \\
        \hline
        Large Language Model based Multi-Agents: A Survey of Progress and Challenges \cite{ref-2} & Significant progress, important challenges, development, effective application. & Complexity of environments, need to improve system efficiency, scalability. & Environment-entity interface, agent profiling, communication, capability acquisition. & Not specified & Not specified & Not specified \\
        \hline
        Optimization of Large Language Model-based Agents \cite{ref-3} & Optimization of LLM agents, improving performance, complex tasks. & Long-term planning, dynamic interaction, adaptability to new scenarios. & Overcome limitations, improve performance, complex tasks. & Not specified & Not specified & Not specified \\
        \hline
        Large Language Model Agent: A Survey on Methodology, Applications and Challenges \cite{ref-4} & Evolving LLM agent research, systematic analysis, architectures, collaboration. & Methodological taxonomy, may limit scope. & Methodological taxonomy, construction, collaboration, evolution of systems. & Not specified & Not specified & Not specified \\
        \hline
        LLM-Based Human-Agent Collaboration and Interaction Systems: A Survey \cite{ref-5} & Lack of research, human-agent collaboration systems, LLMs. & Lack of depth, some aspects of collaboration systems. & Key components, applications, challenges, opportunities. & Not specified & Not specified & Foster research, development of collaboration systems. \\
        \hline
    \end{tabular}
\end{table}

\par\par
Research on Large Language Model (LLM)-based agents has seen growing interest in recent decades, with significant advancements in automating complex tasks. However, despite these efforts, the state of the art presents critical limitations, such as long-term planning, dynamic interaction with environments, and adaptability to new scenarios. LLM agent systems face significant challenges for their effective development and application, including the structural heterogeneity of PDF documents, dependence on the base model's capabilities, and the need to improve efficiency and scalability. In this context, the present work addresses the existing gap in optimizing LLM agents to enhance their performance in complex tasks, providing a systematic review of fine-tuning, reinforcement learning, and other strategies to overcome these limitations.

\begin{thebibliography}{5}
\bibitem{ref-1} Yuheng Cheng, Ceyao Zhang, Zhengwen Zhang, Xiangrui Meng, Sirui Hong, Wenhao Li, Zihao Wang, Zekai Wang, Feng Yin, Junhua Zhao, Xiuqiang He. ``EXPLORING LARGE LANGUAGE MODEL BASED INTELLIGENT AGENTS: D EFINITIONS , METHODS , AND PROSPECTS,'' Not disponible en el texto proporcionado.
\bibitem{ref-2} Taicheng Guo, Xiuying Chen, Yaqi Wang, Ruidi Chang, Shichao Pei, Nitesh V. Chawla, Olaf Wiest. ``Large Language Model based Multi-Agents: A Survey of Progress and Challenges,'' 2023.
\bibitem{ref-3} Shangheng Du, Jiabao Zhao, Jinxin Shi, Zhentao Xie, Xin Jiang, Yanhong Bai, Liang He. ``Optimización de Agentes basados en Lenguaje Grande,'' Nota: El anio no se menciona en el texto..
\bibitem{ref-4} Junyu Luo, Weizhi Zhang, Y e Yuan, Yusheng Zhao, Junwei Yang, Yiyang Gu, Bohan Wu, Binqi Chen, Ziyue Qiao, Qingqing Long, Rongcheng Tu, Xiao Luo, Wei Ju, Zhiping Xiao, Yifan Wang, Meng Xiao, Chenwu Liu, Jingyang Yuan, Shichang Zhang, Yiqiao Jin, Fan Zhang, Xian Wu, Hanqing Zhao, Dacheng Tao. ``Large Language Model Agent: A Survey on Methodology, Applications and Challenges,'' Not especificado.
\bibitem{ref-5} Henry Peng Zou, Wei-Chieh Huang, Yaozu Wu, Yankai Chen, Chunyu Miao, Hoang Nguyen, Yue Zhou, Weizhi Zhang, Liancheng Fang, Langzhou He, Yangning Li, Dongyuan Li, Renhe Jiang, Xue Liu, Philip S. Yu. ``LLM-Based Human-Agent Collaboration and Interaction Systems: A Survey,'' So no se especifica el año de publicación del paper, pero se menciona que fue publicado en junio de 2025.
\end{thebibliography}